\documentclass[20pt]{tikzposter}
\geometry{paperwidth=36in, paperheight=24in}

\usepackage[T1]{fontenc}
\usepackage{lmodern}
\usepackage{amsmath, amssymb}
\usepackage{booktabs}
\usepackage{enumitem}
\usepackage{graphicx}
\usepackage{xcolor}
\usepackage{multicol}

% ── Stanford Colors ─────────────────────────────────────
\definecolor{cardinal}{HTML}{8C1515}
\definecolor{cardinallight}{HTML}{B83A4B}
\definecolor{stanfordwhite}{HTML}{FFFFFF}
\definecolor{stanfordbg}{HTML}{F4F4F4}
\definecolor{stanfordgray}{HTML}{4D4F53}
\definecolor{accent}{HTML}{8C1515}
\definecolor{accentblue}{HTML}{006CB8}
\definecolor{accentorange}{HTML}{E98300}
\definecolor{subtletxt}{HTML}{666666}

\definecolorstyle{Stanford}{
    \definecolor{colorOne}{named}{cardinal}
    \definecolor{colorTwo}{named}{stanfordbg}
    \definecolor{colorThree}{named}{stanfordgray}
}{
    % Background
    \colorlet{backgroundcolor}{stanfordbg}
    \colorlet{framecolor}{cardinal}
    % Title
    \colorlet{titlebgcolor}{cardinal}
    \colorlet{titlefgcolor}{stanfordwhite}
    % Block
    \colorlet{blocktitlebgcolor}{stanfordbg}
    \colorlet{blocktitlefgcolor}{cardinal}
    \colorlet{blockbodybgcolor}{stanfordwhite}
    \colorlet{blockbodyfgcolor}{stanfordgray}
    % Innerblock
    \colorlet{innerblocktitlebgcolor}{stanfordbg}
    \colorlet{innerblocktitlefgcolor}{stanfordgray}
    \colorlet{innerblockbodybgcolor}{stanfordbg}
    \colorlet{innerblockbodyfgcolor}{stanfordgray}
}

\usetheme{Default}
\usecolorstyle{Stanford}
\useblockstyle{Default}

% Compact title area
\definetitlestyle{CompactStanford}{
    width=600mm, roundedcorners=20, linewidth=0.3cm, innersep=0.8cm,
    titletotopverticalspace=4mm, titletoblockverticalspace=6mm,
    titlegraphictotitledistance=10pt, titletextscale=1
}{
    \begin{scope}[line width=\titlelinewidth, rounded corners=\titleroundedcorners]
        \draw[color=framecolor, fill=titlebgcolor]%
        (\titleposleft,\titleposbottom) rectangle (\titleposright,\titlepostop);
    \end{scope}
}
\usetitlestyle{CompactStanford}

% Reduce space between blocks
\makeatletter
\TP@blockverticalspace=6mm
\makeatother

% ── Metadata ────────────────────────────────────────────
\title{\parbox{0.9\linewidth}{\centering
    Exploring Variational Autoencoders for Audio Synthesis and Timbre Transfer
}}
\author{Niccolo Abate \quad\textbar\quad EE\,269 --- CCRMA, Stanford University \quad\textbar\quad March 2026}
\institute{}
\date{}

% ══════════════════════════════════════════════════════════
\begin{document}
\maketitle

\begin{columns}

% ╔══════════════════════════════════╗
% ║  LEFT COLUMN  (0.25)            ║
% ╚══════════════════════════════════╝
\column{0.25}

% ── 1. Motivation ───────────────────────────────────────
\block{1.~Motivation}{
    \textbf{Variational Autoencoders} learn compact latent representations
    enabling reconstruction, synthesis, and \emph{morphing} between sounds ---
    but high-quality waveform decoding remains hard.

    \medskip
    \textbf{Key question:} What is the minimal VAE architecture that yields
    convincing audio \emph{and} a musically useful latent space?

    \medskip
    \textcolor{accentblue}{\textbf{Creative potential:}}
    The latent space as a synthesis interface --- drag a slider to
    continuously morph between any two sounds.

    \medskip
    \textbf{Progression:}
    \begin{itemize}[leftmargin=1.2em, itemsep=2pt]
        \item \textbf{SpecVAE} --- spectrogram-domain baseline
        \item \textbf{RawAudioVAE (small)} --- raw waveform, 29 files
        \item \textbf{RawAudioVAE (ESC-50)} --- scale to 2000 clips, 50 classes
    \end{itemize}
}

% ── 4. Loss Curves ──────────────────────────────────────
\block{4.~Training Loss Curves}{
    \begin{center}
        \includegraphics[width=0.95\linewidth]{rawvae_small/training_curves.png}
        \vspace{1mm}

        {\small\textcolor{subtletxt}{RawAudioVAE (small dataset, 1000 epochs)}}
        \vspace{3mm}

        \includegraphics[width=0.95\linewidth]{specvae/training_curves.png}

        {\small\textcolor{subtletxt}{SpecVAE (myAudioDataset, 232 epochs)}}
    \end{center}

    \medskip
    \textcolor{accentblue}{\textbf{Key:}} $\beta$ warmup (0 $\to$ 0.05 over 100 ep)
    prevents posterior collapse while maintaining reconstruction quality.
    ESC-50 run (not shown) beats small-dataset best in $<$150 epochs.
}

% ── 7. Interpolation ────────────────────────────────────
\block{7.~Latent Space Interpolation}{
    \begin{center}
        \includegraphics[width=0.95\linewidth]{rawvae_small/interp_00_00eac343_to_00ae03f6_grid.png}
    \end{center}

    \medskip
    \textbf{Slerp} between two files' temporal latent sequences
    ($t = 0, 0.25, 0.5, 0.75, 1.0$).
    Each column is a decoded spectrogram at interpolation step $t$.

    \medskip
    \textcolor{accentblue}{\textbf{Mean + offset mode:}} slerp global character
    ($\mu_{\text{global}}$) while blending per-chunk offsets
    $\delta_i = z_i - \mu$ --- smoother transitions, fewer abrupt jumps.
}


% ╔══════════════════════════════════╗
% ║  CENTER COLUMN  (0.50)          ║
% ╚══════════════════════════════════╝
\column{0.50}

% ── 2. Architecture Overview ────────────────────────────
\block{2.~Architecture Overview}{
    \begin{multicols}{2}

    \textcolor{accentblue}{\textbf{SpecVAE}}
    \begin{itemize}[leftmargin=1em, itemsep=1pt]
        \item \textbf{Input:} log-mel spectrogram (64\,bins $\times$ T\,frames)
        \item \textbf{Encoder:} Conv1d stack $\to$ FC $\to$ $\mu, \sigma$
        \item \textbf{Latent:} $z \in \mathbb{R}^{32}$ per chunk (single vector)
        \item \textbf{Decoder:} FC $\to$ DeConv1d $\to$ spectrogram $\to$ Griffin-Lim
        \item \textbf{Loss:} MSE on mel spectrogram + $\beta$-KL
    \end{itemize}

    \columnbreak

    \textcolor{accentorange}{\textbf{RawAudioVAE}}
    \begin{itemize}[leftmargin=1em, itemsep=1pt]
        \item \textbf{Input:} raw waveform ($1 \times 16{,}384$ samples @ 22\,kHz)
        \item \textbf{Encoder:} strided Conv1d ($\times$4, stride 4) + dilated residual
        \item \textbf{Latent:} $z \in \mathbb{R}^{32 \times 64}$ --- \emph{temporal} sequence, $\approx$86\,Hz
        \item \textbf{Decoder:} transposed Conv1d (mirror encoder) $\to$ waveform
        \item \textbf{Loss:} multi-res STFT (2048, 1024, 512, 256) + $\beta$-KL
    \end{itemize}

    \end{multicols}

    \medskip
    \textcolor{accent}{\textbf{Key design choice:}} RawAudioVAE uses a \textbf{temporal latent space} ---
    each latent frame $\approx 11.6$\,ms of audio, preserving fine temporal structure
    that SpecVAE's single-vector latent cannot represent. Free bits ($\lambda = 0.25$) prevent posterior collapse.

    \medskip
    \textit{[Architecture diagrams to be added]}
}

% ── 5. Reconstruction Quality ───────────────────────────
\block{5.~Reconstruction Quality}{
    \begin{center}
    \begin{minipage}{0.48\linewidth}
        \centering
        \includegraphics[width=\linewidth]{specvae/spectrograms.png}\\
        {\small\textcolor{subtletxt}{\textbf{SpecVAE} (ep 232, myAudioDataset)}}
    \end{minipage}\hfill
    \begin{minipage}{0.48\linewidth}
        \centering
        \includegraphics[width=\linewidth]{rawvae_esc50/spectrograms.png}\\
        {\small\textcolor{subtletxt}{\textbf{RawAudioVAE} (ESC-50)}}
    \end{minipage}
    \end{center}

    \medskip
    Each panel: \textbf{Original} (top) $\vert$ \textbf{Reconstructed} (middle) $\vert$ \textbf{Difference} (bottom).
    RawAudioVAE's temporal latent captures fine spectral structure that SpecVAE smooths over.
    Remaining artefacts: low-frequency distortion, addressable with perceptual loss weighting.
}

% ── 8. Demo ─────────────────────────────────────────────
\block{8.~Interactive Demo}{
    \begin{center}
        \textit{[Demo screenshot --- poster/demo\_screenshot.png]}
        \vspace{2mm}

        {\small Upload any audio file, drag the slider to morph between two sounds
        in latent space. Outputs decoded audio + STFT spectrogram + latent PCA scatter.}

        \vspace{3mm}
        \texttt{python demo\_raw.py -r saved/models/RawAudioVAE/<run>/model\_best.pth}

        \vspace{3mm}
        \textit{[QR code to GitHub / demo]}
    \end{center}
}


% ╔══════════════════════════════════╗
% ║  RIGHT COLUMN  (0.25)           ║
% ╚══════════════════════════════════╝
\column{0.25}

% ── 3. Training Setup ───────────────────────────────────
\block{3.~Training Setup}{
    \renewcommand{\arraystretch}{1.25}
    \setlength{\tabcolsep}{4pt}
    \begin{center}
    \small
    \begin{tabular}{l p{1.8cm} p{1.8cm} p{1.8cm}}
        \toprule
        & \textbf{SpecVAE} & \textbf{RawVAE small} & \textbf{RawVAE ESC-50} \\
        \midrule
        Dataset   & myAudio & myAudio & ESC-50 \\
        Files     & 29 & 29 & 2{,}000 \\
        Classes   & 2  & 2  & 50 \\
        Chunks    & $\sim$165 & $\sim$165 & $\sim$9{,}600 \\
        Latent    & $\mathbb{R}^{32}$ & $\mathbb{R}^{32\times64}$ & $\mathbb{R}^{32\times64}$ \\
        Recon loss & MSE (mel) & Multi-res & Multi-res \\
                   &           & STFT & STFT \\
        $\beta$ warmup & --- & $0{\to}0.05$ & $0{\to}0.05$ \\
                       &     & /100\,ep & /100\,ep \\
        Free bits & --- & $\lambda{=}0.25$ & $\lambda{=}0.25$ \\
        Best val  & --- & 1.497 & 1.355 \\
        \quad recon &   & (ep\,675) & (ep\,143+) \\
        \bottomrule
    \end{tabular}
    \end{center}

    \medskip
    \textcolor{accentblue}{\textbf{Note:}} Free bits enforce a minimum KL per latent
    dimension, preventing dimensions from collapsing to the prior.
    ESC-50 generalises far better --- beats 1000-epoch small-data run in $<$150 epochs.
}

% ── 6. Latent Space Structure ───────────────────────────
\block{6.~Latent Space Structure}{
    \begin{center}
        \includegraphics[width=0.95\linewidth]{rawvae_esc50/kl_per_dim.png}
        {\small\textcolor{subtletxt}{KL per latent dim --- ESC-50 ep 100}}
        \vspace{3mm}

        \includegraphics[width=0.95\linewidth]{rawvae_esc50/latent_pca.png}
        {\small\textcolor{subtletxt}{Latent PCA scatter (colour = class, unsupervised)}}
    \end{center}

    \medskip
    Free bits enforce a floor at $\lambda{=}0.25$ nats per dim.
    PCA scatter shows \textbf{class separation emerges unsupervised}
    --- no class labels used during training.
}

% ── 9. Conclusions & Future Work ────────────────────────
\block{9.~Conclusions \& Future Work}{
    \textcolor{accentblue}{\textbf{What worked:}}
    \begin{itemize}[leftmargin=1.2em, itemsep=1pt]
        \item Raw waveform VAE with multi-resolution STFT loss trains stably
        \item Temporal latent space richer than single-vector SpecVAE
        \item ESC-50 generalises far better; beats 1000-ep small-data run in $<$150\,ep
        \item Free bits + $\beta$ warmup prevents posterior collapse
    \end{itemize}

    \medskip
    \textcolor{accentorange}{\textbf{Future work:}}
    \begin{itemize}[leftmargin=1.2em, itemsep=1pt]
        \item Adversarial / feature-matching loss (RAVE-style) for perceptual quality
        \item Perceptual weighting to address low-frequency distortion
        \item Larger latent dim (64+) for diverse datasets
        \item Streaming / real-time inference
    \end{itemize}

    \medskip
    \textcolor{subtletxt}{\small
    \textbf{References:} Kingma \& Welling, VAE (2013);
    Caillon \& Esling, RAVE (2021);
    Défossez et al., EnCodec (2022).
    }
}

\end{columns}

\end{document}
